\chapter{态射与正则性的补充}\label{chap:ag-3}
\setcounter{footnote}{2}
% 原PDF p35.
\textbf{约定：}为简便起见，以下出现的所有概形均假设为局部 Noether 概形，所有态射均假设为分离且有限型的态射，或局部态射.
\section{余法层、正则浸入与正则概形}\label{sec:ag-3-1}
\begin{numpara}[$I$-进滤过]\label{item:ag-3-1-1}
先回顾一些交换代数知识. 设 $A$ 是环，$I$ 是理想. 任意 $A$-模 $M$ 都可赋予（递降的）\newterm{ag-i-adic-filtration}{$I$-进滤过}{$I$-adic filtration} $(I^nM)_{n\in\mathbb N}$：
\begin{align*}
M\supset IM\supset I^2M\supset\cdots.
\end{align*}
当 $A=M$ 时，这是环 $A$ 的滤过. 回顾这些滤过的两条性质：
\begin{enumerate}[label=\roman*)]
\item \emph{Artin--Rees 引理：}设 $A$ 为 Noether 环，$M$ 有限生成，$N$ 是 $M$ 的子模. 则存在整数 $n_0$，使得对所有 $n\ge n_0$，有 $I^nM\cap N=I^{n-n_0}(I^{n_0}M\cap N)$.
\item \emph{Krull 交定理：}设 $A$ 为 Noether 环，$M$ 有限生成，且 $I$ 包含于 $A$ 的所有极大理想. 则 $M$ 上的 $I$-进滤过是分离的，即 $\bigcap_{n\in\mathbb N}I^nM=\{0\}$.（这可由 Artin--Rees 引理与 Nakayama 引理推出.）
\end{enumerate}
现在定义：
\begin{itemize}
\item \newterm{ag-associated-graded-module}{结合分次模}{associated graded module}
\begin{align*}
\operatorname{gr}^{I}(M)&:=M/IM\oplus IM/I^2M\oplus I^2M/I^3M\oplus\cdots\\
&=\bigoplus_{n\in\mathbb N}I^nM/I^{n+1}M.
\end{align*}
它是分次 $A/I$-模；当 $M=A$ 时，它是分次 $A/I$-代数.
\item 伴随的\newterm{ag-rees-module}{Rees 模}{Rees module}
\begin{align*}
\operatorname{Rees}^{I}(M):=M\oplus IM\oplus I^2M\oplus\cdots=\bigoplus_{n\in\mathbb N}I^nM.
\end{align*}
它是分次 $A$-模；当 $M=A$ 时，它就是前面已经遇到的分次 $A$-代数.
\end{itemize}
显然，存在分次 $A/I$-模的典范同构
\begin{align}
A/I\otimes_A\operatorname{Rees}^{I}(M)\xrightarrow\sim\operatorname{gr}^{I}(M),\tag{3.1.1.1}\label{eq:ag-3-1-1-1}
\end{align}
当 $M=A$ 时，它是分次 $A/I$-代数的同构.
对称代数的万有性质给出两个态射，分别为分次 $A$-代数态射与分次 $A/I$-代数态射：
\begin{align}
\operatorname{Sym}_A(I)\longrightarrow\operatorname{Rees}^{I}(A)
\quad\text{与}\quad
\operatorname{Sym}_{A/I}(I/I^2)\longrightarrow\operatorname{gr}^{I}(A).\tag{3.1.1.2}\label{eq:ag-3-1-1-2}
\end{align}
其中第二个由第一个施行张量运算 $A/I\otimes-$ 得到. 这些态射显然满射.
这些构造与局部化相容，更一般地，与平坦基变换相容. 事实上，对任意在 $A$ 上平坦的 $A$-代数 $A'$ 及任意 $n\in\mathbb N$，记 $I':=IA$ 为 $I$ 在 $A'$ 中生成的理想，并记 $M':=A'\otimes_AM$，则有典范同构 $A'\otimes_A(I^nM)\xrightarrow\sim I'^nM'$ 以及
\begin{align}
(A'/I')\otimes_{A/I}(I^nM/I^{n+1}M)\xrightarrow\sim I'^nM'/I'^{n+1}M'.\tag{3.1.1.3}\label{eq:ag-3-1-1-3}
\end{align}
% 原PDF p36.
由于与局部化相容，可将上述构造如下层化. 设 $i:Y\hookrightarrow X$ 是概形的闭浸入，$\mathcal I:=\operatorname{Ker}(\mathcal O_X\longrightarrow i_*\mathcal O_Y)$ 是 $X$ 上对应的 $\mathcal O_X$-模层. 任意 $\mathcal O_X$-模 $\mathcal M$ 均带有 $\mathcal I$-进滤过 $(\mathcal I^n\mathcal M)_{n\in\mathbb N}$，其结合分次 $\mathcal O_X$-模为 $\operatorname{gr}^{\mathcal I}(\mathcal M):=\bigoplus_n\mathcal I^n\mathcal M/\mathcal I^{n+1}\mathcal M$，Rees $\mathcal O_X$-模为 $\operatorname{Rees}^{\mathcal I}(\mathcal M):=\bigoplus_{n\in\mathbb N}\mathcal I^n\mathcal M$. 当 $\mathcal M$ 在 $X$ 上拟凝聚时，这些层也全都拟凝聚. 特别地，若 $U=\operatorname{Spec}(A)$ 是仿射开集，$I=\mathcal I(U)$ 是给出闭浸入 $U\cap Y\hookrightarrow U$ 的理想，则 $\operatorname{gr}^{\mathcal I}(\mathcal M)(U)=\operatorname{gr}^{I}(\mathcal M(U))$，等等. 各仿射坐标图上的同构~\eqref{eq:ag-3-1-1-1}于是诱导同构
\begin{align}
i^*(\operatorname{Rees}^{\mathcal I}(\mathcal M))\xrightarrow\sim i^{-1}(\operatorname{gr}^{\mathcal I}(\mathcal M)).\tag{3.1.1.4}\label{eq:ag-3-1-1-4}
\end{align}
同样，满态射~\eqref{eq:ag-3-1-1-2}诱导满态射
\begin{align}
\begin{aligned}
\text{(i)}\quad&\operatorname{Sym}_{\mathcal O_X}(\mathcal I)\longrightarrow\operatorname{Rees}^{\mathcal I}(\mathcal O_X),\\
\text{(ii)}\quad&\operatorname{Sym}_{\mathcal O_Y}(i^{-1}(\mathcal I/\mathcal I^2))\longrightarrow i^{-1}(\operatorname{gr}^{\mathcal I}(\mathcal O_X)).
\end{aligned}\tag{3.1.1.5}\label{eq:ag-3-1-1-5}
\end{align}
以后有时会省略 $i^{-1}$，这也是文献中常见的记号滥用.
\emph{$\operatorname{gr}^{\mathcal I}(\mathcal O_X)$ 的几何解释.} 设 $\widetilde X_Y\xrightarrow{\pi_Y}X$ 是 $X$ 沿 $Y$ 的爆破. 上述同构 (ii) 将例外除子 $\pi_Y^{-1}(Y)\longrightarrow Y$ 等同于 $\operatorname{Proj}(\operatorname{gr}^{\mathcal I}(\mathcal O_X))$.
\end{numpara}
\begin{definition}\label{item:ag-3-1-2}
设 $i:Y\hookrightarrow X$ 是闭浸入，$\mathcal I:=\operatorname{Ker}(\mathcal O_X\longrightarrow i_*\mathcal O_Y)$ 是 $X$ 上对应的 $\mathcal O_X$-模层. 称 $Y$ 上的拟凝聚 $\mathcal O_Y$-模
\begin{align*}
\mathcal N_i=\mathcal N_{Y/X}:=i^*(\mathcal I)=i^{-1}(\mathcal I/\mathcal I^2)
\end{align*}
为 $Y$ 在 $X$ 中的\newterm{ag-conormal-sheaf}{余法层}{conormal sheaf}\footnote{这里采用 EGA IV 的记号. 在 Hartshorne 的书中，此记号表示法层（余法层的对偶）.}.
\end{definition}
\begin{example}\label{item:ag-3-1-3}
\begin{enumerate}[label=(\roman*)]
\item 设 $Y=\{x\}$ 是 $X$ 的一个闭点. 则 $\mathcal N_{Y/X}=\mathfrak m_x/\mathfrak m_x^2$ 是 $X$ 在 $x$ 处的余切空间（一个 $\kappa(x)$-向量空间）.
\item 设 $k$ 是域，且 $Y=\mathbb A_k^m=\operatorname{Spec}(k[T_1,\ldots,T_m])\subset\mathbb A_k^n=\operatorname{Spec}(k[T_1,\ldots,T_n])=X$. 则可验证，$k[T_1,\ldots,T_m]$-线性映射
\begin{align*}
k[T_1,\ldots,T_m]^{n-m}&\longrightarrow\mathcal N_{Y/X}=(T_{m+1},\ldots,T_n)/(T_{m+1},\ldots,T_n)^2,\\
(P_{m+1},\ldots,P_n)&\longmapsto\sum_{i=m+1}^nT_iP_i
\end{align*}
是同构. 因而这里 $\mathcal N_{Y/X}$ 是秩为 $n-m$ 的自由模.
\item 设 $Y$ 是 $X$ 中的一个除子. 此时 $\mathcal N_{Y/X}$ 是 $Y$ 上的可逆层，因为理想 $\mathcal I$ 本身在 $X$ 上可逆.\unskip\nobreak\hfill\ensuremath{\diamond}
\end{enumerate}
\renewcommand{\envenddiamond}{}
\end{example}
% 原PDF p37.
\begin{numpara}[余法层的函子性]\label{item:ag-3-1-4}
由于我们的概形通常假设为 Noether 概形，$\mathcal N_{Y/X}$ 是 $Y$ 上的凝聚层. 注意，若
\begin{align}
\begin{tikzcd}[ampersand replacement=\&]
Y'\arrow[r,hook,"i'"]\arrow[d,"g"'] \& X'\arrow[d,"f"]\\
Y\arrow[r,hook,"i"'] \& X
\end{tikzcd}\tag{3.1.4.1}\label{eq:ag-3-1-4-1}
\end{align}
是交换图，则有典范态射
\begin{align}
g^*\mathcal N_{Y/X}\longrightarrow\mathcal N_{Y'/X'}.\tag{3.1.4.2}\label{eq:ag-3-1-4-2}
\end{align}
事实上，只须在仿射情形构造它；采用显然的记号，它由态射 $A'\otimes_AI\longrightarrow I'$ 诱导. 当此图是笛卡尔图时，在仿射情形即要求 $I'=IA'$，此态射显然满射. 以下两种重要情形中，它甚至是同构.
\end{numpara}
\begin{proposition}\label{item:ag-3-1-5}
沿用上述记号，假设处于下述两种情形之一：
\begin{enumerate}[label=\roman*)]
\item 图~\eqref{eq:ag-3-1-4-1}是笛卡尔图，且 $f$ 平坦.
\item 图~\eqref{eq:ag-3-1-4-1}嵌入图
\begin{align*}
\begin{tikzcd}[ampersand replacement=\&]
Y'\arrow[r,hook,"i'"]\arrow[d,"g"'] \& X'\arrow[r,"p'"]\arrow[d,"f"] \& Y'\arrow[d,"g"]\\
Y\arrow[r,hook,"i"'] \& X\arrow[r,"p"'] \& Y
\end{tikzcd}
\end{align*}
其中 $p\circ i=\operatorname{Id}_Y$，且两个方块都是笛卡尔方块.
\end{enumerate}
则态射~\eqref{eq:ag-3-1-4-2}是同构.
\end{proposition}
\begin{proof}
两种情形都只须处理仿射情形. 情形 i) 是~\eqref{eq:ag-3-1-1-3}的特例. 在情形 ii) 中，设 $Y=\operatorname{Spec}(B)$，$X=\operatorname{Spec}(A)$，$I$ 是 $A$ 的理想，使得 $B\longrightarrow A\longrightarrow A/I$ 是同构. 于是 $A=B\oplus I$. 再设 $Y'=\operatorname{Spec}(B')$，则 $A'=B'\oplus(B'\otimes_BI)$. 特别地，满态射 $A'\otimes_AI=B'\otimes_BI\longrightarrow I'=IA'$ 是同构，从而 $A'\otimes_A(I/I^2)\xrightarrow\sim I'/I'^2$. 不过，若没有额外假设，典范态射 $A'\otimes_A\operatorname{gr}^{I}(A)\longrightarrow\operatorname{gr}^{I'}(A')$ 没有理由是双射.
\end{proof}
\begin{definition}\label{item:ag-3-1-6}
若 $\mathcal O_Y$-模 $\mathcal N_{Y/X}$ 局部自由（且秩有限），并且态射~\eqref{eq:ag-3-1-1-5}(ii) 是同构，则称浸入 $i:Y\hookrightarrow X$ 为\newterm{ag-regular-immersion}{正则浸入}{regular immersion}.
\end{definition}
\noindent\textbf{例：}若 $Y$ 是除子，则 $i$ 为正则浸入.
下面的命题描述正则浸入的闭子概形的爆破的例外除子. 这推广了前面在仿射空间一个点的爆破的例子中所作的观察.
% 原PDF p38.
\begin{proposition}\label{item:ag-3-1-7}
设 $i:Y\hookrightarrow X$ 是正则闭浸入，$Y'$ 是 $\widetilde X_Y$ 的例外除子（即 $Y$ 的逆像）.
\begin{enumerate}[label=\roman*)]
\item 存在 $Y$-概形的典范同构 $Y'\xrightarrow\sim\mathbb P(\mathcal N_{Y/X})$.
\item 在此等同之下，可逆层 $\mathcal O_{\mathbb P(\mathcal N_{Y/X})}(1)$ 对应于 $\mathcal N_{Y'/\widetilde X_Y}$.
\end{enumerate}
\end{proposition}
\begin{proof}
设 $\mathcal I$ 是 $Y$ 的理想. 有
\begin{align*}
Y'&=i^*\operatorname{Proj}(\operatorname{Rees}^{\mathcal I}(\mathcal O_X))
=\operatorname{Proj}(i^*\operatorname{Rees}^{\mathcal I}(\mathcal O_X))\\
&=\operatorname{Proj}(i^{-1}\operatorname{gr}^{\mathcal I}(\mathcal O_X)).
\end{align*}
由于假设 $i$ 正则，$\mathcal O_Y$-代数态射
\begin{align*}
\operatorname{Sym}_{\mathcal O_Y}(i^*\mathcal I)\longrightarrow i^{-1}\operatorname{gr}^{\mathcal I}(\mathcal O_X)
\end{align*}
是同构. 因 $\mathcal N_{Y/X}=i^*(\mathcal I)$，第一条断言随即得出. 除子 $Y'$ 由分次理想 $\mathcal I'=\mathcal I\operatorname{Rees}^{\mathcal I}(\mathcal O_X)=\bigoplus_n\mathcal I^{n+1}$ 所伴随的理想层定义，此分次理想是 $\operatorname{Rees}^{\mathcal I}(\mathcal O_X)$ 的理想. 因此，$i'^*(\mathcal I')$ 是 $i^{-1}\operatorname{gr}^{\mathcal I}(\mathcal O_X)$ 上的分次模 $\bigoplus_n\mathcal I^{n+1}/\mathcal I^{n+2}$ 所伴随的模层. 此分次模到 $\operatorname{gr}^{\mathcal I}(\mathcal O_X)(1)$ 的分次模态射诱导同构 $i'^*(\mathcal I')\xrightarrow\sim\mathcal O_{\mathbb P(i^*\mathcal I)}(1)$.
\end{proof}
回到正则浸入的定义. 注意，所要求的条件具有局部性. 更确切地说，给定 $Y$ 的一点 $x$，以下两个性质等价：
\begin{enumerate}[label=\roman*)]
\item 存在 $x$ 在 $X$ 中的开邻域 $U$，使浸入 $Y\cap U\hookrightarrow U$ 正则.
\item 闭浸入 $\operatorname{Spec}(\mathcal O_{Y,x})\hookrightarrow\operatorname{Spec}(\mathcal O_{X,x})$ 正则.
\end{enumerate}
此时称浸入 $Y\hookrightarrow X$ \emph{在点 $x$ 处正则}. 立即可见，浸入正则当且仅当它在每一点处正则. 因而只须理解 $X=\operatorname{Spec}(A)$、$Y=\operatorname{Spec}(A/I)$ 且 $A$ 为局部环的情形. 此时有如下结论.
\begin{proposition}\label{item:ag-3-1-8}
设 $A$ 是 Noether 局部环，$I$ 是不等于 $A$ 的理想. 浸入 $\operatorname{Spec}(A/I)\hookrightarrow\operatorname{Spec}(A)$ 正则（也称 $I$ 是 $A$ 的\newterm{ag-regular-ideal}{正则理想}{regular ideal}），当且仅当 $I$ 存在一族生成元 $f_1,\ldots,f_m$，满足下列等价性质之一：
\begin{enumerate}[label=\roman*)]
\item 由 $f_i$ 定义的典范态射 $(A/I)[T_1,\ldots,T_m]\longrightarrow\operatorname{gr}^{I}(A)$ 是同构.（这里 $T_i$ 为不定元.）
\item 序列 $(f_1,\ldots,f_m)$ 正则，参见分目~\ref{item:ag-2-2-4}.
\end{enumerate}
若还给定 $A$ 的子代数 $B$，使 $B\longrightarrow A\longrightarrow A/I$ 是同构，则上述条件又等价于
\begin{enumerate}[label=\roman*),start=3]
\item 由 $f_i$ 定义的典范态射 $B[[T_1,\ldots,T_m]]\longrightarrow\varprojlim_n(A/I^n)$ 是同构.
\end{enumerate}
在这些条件下，有 $\dim(A/I)=\dim(A)-m$.
\end{proposition}
注意等价关系 i)$\Leftrightarrow$ii) 的如下推论：在 Noether 局部环中，一列元素是否正则与这些元素的顺序无关.
% 原PDF p39.
\begin{proof}
浸入 $\operatorname{Spec}(A/I)\hookrightarrow\operatorname{Spec}(A)$ 的正则性显然等价于 $I$ 中存在一族 $f_1,\ldots,f_n$ 满足条件 i)（取 $f_i$ 使它们在 $I/I^2$ 中的像构成一组基）.
另一方面，给定这样一族元素，蕴含 ii)$\Rightarrow$i) 由分目~\ref{item:ag-2-2-4}中的命题得到.
现在证明 i)$\Rightarrow$ii). 为此，如分目~\ref{item:ag-2-2-4}中的命题那样作归纳. 令 $J:=(f_1,\ldots,f_{m-1})$，考虑分解
\begin{align*}
\begin{tikzcd}[ampersand replacement=\&,column sep=large]
A[T_1,\ldots,T_m]\arrow[r] \& \operatorname{Rees}^{J}(A)[T_m]\arrow[r,"\alpha"] \& \operatorname{Rees}^{I}(A)
\end{tikzcd}
\end{align*}
其中 $\alpha$ 将 $T_m$ 映到 $f_mt$，并延拓 $\operatorname{Rees}^{J}(A)$ 到 $\operatorname{Rees}^{I}(A)$ 的典范包含. 与 $A/I$ 作张量积后，得到交换图
\begin{align*}
\begin{tikzcd}[ampersand replacement=\&,column sep=large]
A/J[T_1,\ldots,T_{m-1}]\otimes_A A/(f_m)[T_m]\arrow[r,"\psi_J\otimes\operatorname{Id}"]\arrow[d,equal] \& \operatorname{gr}^{J}(A)\otimes_A A/(f_m)[T_m]\arrow[d,"\beta"]\\
A/I[T_1,\ldots,T_m]\arrow[r,"\psi_I"'] \& \operatorname{gr}^{I}(A)
\end{tikzcd}
\end{align*}
由假设，$\psi_I$ 是同构. 此图表明 $\psi_J$ 单射，从而是同构. 归纳假设于是说明 $(f_1,\ldots,f_{m-1})$ 是正则序列. 还须证明 $f_m$ 在 $A/J$ 中不是零因子（即乘以 $f_m$ 在 $A/J$ 上是单射）. 为此，注意 $\beta$ 是双射. 简记 $f:=f_m$、$T:=T_m$，取 $a\in A$ 使 $af\in J$，写成 $af+j=0$. 这蕴含 $\operatorname{gr}^{J}(A)\otimes_AA/(f)[T]$ 中一次元素 $aT\oplus j$ 在 $\operatorname{gr}^{I}(A)$ 中的像为零；由 $\beta$ 的单射性，得到 $a\in I=(f)+J$ 以及 $j\in JI$. 写成 $a=a'f+j'$. 此时 $\operatorname{gr}^{J}(A)\otimes_AA/(f)[T]$ 中的二次元素 $a'T^2\oplus j'T\oplus j$ 在 $\operatorname{gr}^{I}(A)$ 中的像为零；由 $\beta$ 的单射性，得到 $a'\in I=(f)+J$，从而 $a\in(f^2)+J$. 继续这一逐次分解过程，归纳可得对所有 $n\in\mathbb N$，均有 $a\in(f^n)+J$. 而环 $A/J$ 是 Noether 环，所以由上面回顾的 Krull 定理，$f$-进滤过是分离的（即 $\bigcap_{n\in\mathbb N}f^nA/J=\{0\}$），故 $a\in J$. 因此 $f$ 在 $A/J$ 中不是零因子.
在适当假设下，ii) 与 iii) 的等价性很容易，留给读者证明.
最后，维数等式来自如下事实：若 $f\in A$ 不是零因子，则 $\dim(A/(f))=\dim(A)-1$.
\end{proof}
% 原PDF p39–44：§3.1.9–3.1.15.
在 $I$ 是 $A$ 的极大理想 $\mathfrak m$ 这一特殊情形下，还可得到更多结论：
\begin{proposition}\label{item:ag-3-1-9}
设 $A$ 是极大理想为 $\mathfrak m$ 的 Noether 局部环，则下列条件等价：
\begin{enumerate}[label=\roman*)]
\item 理想 $\mathfrak m$ 正则.
\item $\dim_{A/\mathfrak m}(\mathfrak m/\mathfrak m^2)=\dim(A)$（即局部环 $A$ 是\newterm{ag-regular-local-ring}{正则局部环}{regular local ring}）.
\end{enumerate}
\end{proposition}
\begin{proof}
令 $d=\dim_{A/\mathfrak m}(\mathfrak m/\mathfrak m^2)$，取 $\mathfrak m$ 中的元素 $f_1,\ldots,f_d$，使其像构成 $\mathfrak m/\mathfrak m^2$ 的一组 $A/\mathfrak m$-基. 由 Nakayama 引理，它们生成 $\mathfrak m$. 考虑相应的分次满态射
\begin{align}
A/\mathfrak m[T_1,\ldots,T_d]\longrightarrow\operatorname{gr}^{\mathfrak m}(A).\tag{3.1.9.3}\label{eq:ag-3-1-9-3}
\end{align}
由 Hilbert--Samuel 多项式理论（参见 Eisenbud，第 12 章），存在次数为 $\dim(A)-1$ 的多项式 $P$，使得对充分大的整数 $n$，有 $\dim(\mathfrak m^n/\mathfrak m^{n+1})=P(n)$.
假设 $\mathfrak m$ 正则. 则~\eqref{eq:ag-3-1-9-3}是同构，故 $P(n)=\binom{n+d-1}{d-1}$，其次数为 $d-1$，从而 $d=\dim(A)$. 反过来，假设 $d=\dim(A)$. 取~\eqref{eq:ag-3-1-9-3}的核中一个非零齐次元 $a$，设其次数为 $m$. 记多项式代数 $A/\mathfrak m[T_1,\ldots,T_d]$ 的分次为 $\bigoplus_{k\in\mathbb N}S_k$. 当 $n>m$ 时，$S_n$ 的子模 $aS_{n-m}$ 包含在~\eqref{eq:ag-3-1-9-3}的核中. 若 $a$ 非零，则 $\dim(aS_{n-m})=\dim(S_{n-m})=\binom{n+d-1-m}{d-1}$. 因而有不等式
\begin{align*}
\dim(\mathfrak m^n/\mathfrak m^{n+1})\leqslant\binom{n+d-1}{d-1}-\binom{n+d-1-m}{d-1}.
\end{align*}
但右边是关于 $n$ 的次数小于 $d-1$ 的多项式，与 $d=\dim(A)$ 矛盾. 所以 $a$ 为零，故~\eqref{eq:ag-3-1-9-3}确为单射.
\end{proof}
现在列举（回顾）正则局部环的一些性质. 若 $\mathfrak m$ 中元素所成正则序列的最大长度 $k$（称为 $A$ 的\newterm{ag-depth}{深度}{depth}）等于 $\dim(A)$，则称局部环 $A$ 为\newterm{ag-cohen-macaulay-ring}{Cohen--Macaulay 环}{Cohen--Macaulay ring}. 我们将用到如下事实：若 $f\in A$ 不是零因子，则 $\dim(A/(f))=\dim(A)-1$.
\begin{theorem}\label{item:ag-3-1-10}
设 $A$ 是 Noether 正则局部环，则：
\begin{enumerate}[label=\roman*)]
\item $A$ 是整环、整闭环（正规环），且是 Cohen--Macaulay 环.
\item 若 $I$ 是 $A$ 的理想，则 $A/I$ 正则当且仅当 $I$ 有一组生成元 $(f_1,\ldots,f_k)$，它们模 $\mathfrak m^2$ 的类在 $A/\mathfrak m$ 上线性无关. （这是 Jacobi 判别法的抽象形式，参见下例.）
\item 若 $A=B/J$，其中 $B$ 是 Noether 局部环，$J$ 是理想，则 $B$ 正则当且仅当 $J$ 是 $B$ 的正则理想.
\item $A$ 的上同调维数有限. （此性质刻画了正则环，且还有 $\operatorname{dim.coh.}(A)=\dim(A)$.）
\item 对 $A$ 的每个素理想 $\mathfrak p$，局部环 $A_{\mathfrak p}$ 都正则.
\item $A[T_1,\ldots,T_k]$ 的每个局部化都正则.
\end{enumerate}
\end{theorem}
\begin{proof}
\textup{i)} 一般地，若一个环具有分离滤过，且其结合分次环是整环（分别地，整闭环），则此环也是整环（分别地，整闭环），参见 Bourbaki，\emph{Alg. Com.}，第 III 章第 2 节命题 1 的推论，以及第 V 章第 1 节第 5 目命题 15. 此处，因 $A$ 是 Noether 环，$\mathfrak m$-进滤过分离（见前面回顾的 Krull 定理），而结合分次环 $\operatorname{gr}^{\mathfrak m}(A)$ 是多项式代数，因而是整环且正规. 此外，由定理前回顾的性质可知，若 $(f_1,\ldots,f_k)$ 是正则序列，则 $\dim(A/(f_1,\ldots,f_k))=\dim(A)-k$；特别地，这说明 $\operatorname{prof}(A)\leqslant\dim(A)$. 但由命题~\ref{item:ag-3-1-9}，存在长度为 $\dim(A)=\dim_{A/\mathfrak m}(\mathfrak m/\mathfrak m^2)$ 的正则序列. 所以 $A$ 是 Cohen--Macaulay 环.
\textup{ii)} 有正合列
\begin{align*}
(\mathfrak m^2+I)/\mathfrak m^2\hookrightarrow\mathfrak m/\mathfrak m^2\twoheadrightarrow(\mathfrak m/I)/(\mathfrak m/I)^2.
\end{align*}
假设 $I$ 由陈述中的一族元素 $(f_1,\ldots,f_k)$ 生成，则有 $\dim((\mathfrak m/I)/(\mathfrak m/I)^2)=\dim(\mathfrak m/\mathfrak m^2)-k=\dim(A)-k$. 又因理想 $\mathfrak m$ 正则，由命题~\ref{item:ag-3-1-8}可知序列 $(f_1,\ldots,f_k)$ 正则，从而 $\dim(A/I)=\dim(A)-k$. 所以 $A/I$ 正则. 反过来，假设 $A/I$ 正则，令 $k=\dim(A)-\dim(A/I)$. 由上述正合列，$k=\dim_{A/\mathfrak m}(\mathfrak m^2+I/\mathfrak m^2)$. 取 $I$ 中元素 $(f_1,\ldots,f_k)$，使其提升 $(\mathfrak m^2+I)/\mathfrak m^2$ 的一组基. 它们构成正则序列，因而由刚才的证明，它们生成理想 $I'\subset I$，使得 $A/I'$ 正则且维数为 $d-k$. 但 $A/I'$ 是整环，并且与 $A/I$ 维数相同，故投影 $A/I'\twoheadrightarrow A/I$ 是同构，所以 $I=I'$.
\textup{iii)} 方向 $\Rightarrow$ 由 ii) 得到. 现假设 $J$ 正则.\footnote{这里有一个学生在课堂上找到的更简单的论证. 选取生成 $J$ 的正则序列，再用 $A$ 中生成 $\mathfrak m$ 的正则序列的提升将其补全. 所得序列正则且生成 $B$ 的极大理想，故由前一命题，此极大理想正则.}于是 $J$ 由正则序列 $(f_1,\ldots,f_k)$ 生成，且有 $\dim(A)+k=\dim(B)$. 设 $\mathfrak n$ 为 $B$ 的极大理想. 有正合列
\begin{align*}
(\mathfrak n^2+J)/\mathfrak n^2\hookrightarrow\mathfrak n/\mathfrak n^2\twoheadrightarrow\mathfrak m/\mathfrak m^2.
\end{align*}
有 $\dim_{B/\mathfrak n}((\mathfrak n^2+J)/\mathfrak n^2)\leqslant\operatorname{rg}_{B/J}(J/J^2)=k$，而因 $A$ 正则，$\dim(\mathfrak m/\mathfrak m^2)=\dim(A)$. 由此得到 $\dim(\mathfrak n/\mathfrak n^2)\leqslant\dim(B)$. 另一方面，一般地总有不等式 $\dim(\mathfrak n/\mathfrak n^2)\geqslant\dim(B)$（它由 Hilbert 多项式理论得到），所以 $B$ 正则.
性质 iv)、v)、vi) 的特点是需要使用上同调方法. 注意，iv)$\Rightarrow$v)，因为总有 $\operatorname{dim.coh.}(A_{\mathfrak p})\leqslant\operatorname{dim.coh.}(A)$. 同样，iv)$\Rightarrow$vi)，因为总有 $\operatorname{dim.coh.}(A[t])\leqslant\operatorname{dim.coh.}(A)+1$.
iv) 的证明的核心，是与 $A$ 中一族元素 $(f_1,\ldots,f_k)$ 相伴的\newterm{ag-koszul-complex}{Koszul 复形}{Koszul complex} $K^\bullet(f_1,\ldots,f_k)$. 它定义为复形 $K^\bullet(f_i):=[A\xrightarrow{\times f_i}A]$ 的张量积，各复形置于次数 $-1,0$. 容易验证，若此序列正则，则该复形的上同调集中于次数 $0$，且为 $A/(f_1,\ldots,f_k)$. 在此情形下，便得到 $A$-模 $A/(f_1,\ldots,f_k)$ 的长度为 $k$ 的投射解消. 此外，复形 $K^\bullet(f_1,\ldots,f_k)\otimes A/(f_1,\ldots,f_k)$ 的微分为零，所以
\begin{align*}
\operatorname{Tor}_A^k(A/(f_1,\ldots,f_k),A/(f_1,\ldots,f_k))\neq0,
\end{align*}
因而 $\operatorname{dim.proj.}(A/(f_1,\ldots,f_k))=k$. 另一方面，可用逐步分解的方法证明，对每个 Noether 局部环，都有 $\operatorname{dim.coh.}(A)=\operatorname{proj.dim.}_A(A/\mathfrak m)$. 因此，若 $A$ 正则，则其上同调维数有限，且恰等于 $\dim(A)$. 逆命题更为困难. 首先证明，若 $\operatorname{dim.coh.}(A)$ 有限，则它不大于 $\dim(A)$. 然后选取 $\mathfrak m$ 中一族元素 $(f_1,\ldots,f_k)$，使它们诱导 $\mathfrak m/\mathfrak m^2$ 的一组基（因而 $k\geqslant\dim(A)$）. Koszul 复形 $K^\bullet(f_1,\ldots,f_k)$ 未必仍是 $A$-模 $A/\mathfrak m$ 的解消，但通过研究相应的谱序列，可以证明
\begin{align*}
\operatorname{Tor}^k(A/\mathfrak m,A/\mathfrak m)\neq0,
\end{align*}
从而 $\operatorname{dim.coh.}(A)\geqslant k$. 这只有在 $k=\dim(A)$ 时才可能成立，即 $A$ 正则.
\end{proof}
\begin{remark*}
正则性的上同调刻画，以及更一般地，深度的上同调解释，都是非常重要的：它们是强有力的工具，而且使这些概念可以推广到非交换 Noether 环.
\end{remark*}
\begin{remark}\label{item:ag-3-1-11}
上述略显难懂的交换代数证明中使用或提及的技术，最应记住的是存在下列对象和方法：
\begin{itemize}
\item 分次模的 Hilbert 多项式，或局部环的 Hilbert--Samuel 多项式.
\item 将结合分次环的某些性质转移到带有``良好''滤过的环上的方法.
\item 与有限族元素相伴的 Koszul 复形.\unskip\nobreak\hfill\ensuremath{\diamond}
\end{itemize}
\renewcommand{\envenddiamond}{}
\end{remark}
\begin{corollary}\label{item:ag-3-1-12}
设 $X$ 是 Noether 概形.
\begin{enumerate}[label=\roman*)]
\item $X$ 正则当且仅当对 $X$ 的每个闭点 $x$，闭浸入 $\operatorname{Spec}(\kappa(x))\hookrightarrow X$ 正则.
\item 设 $i:Y\hookrightarrow X$ 是闭浸入，$x$ 是 $Y$ 的正则点. 则 $i$ 在 $x$ 处正则当且仅当 $X$ 在 $x$ 处正则.
\end{enumerate}
\end{corollary}
\begin{proof}
\textup{i)} 回顾，若 $X$ 的所有局部环都正则，则称 $X$ 为\newterm{ag-regular-scheme}{正则概形}{regular scheme}. 因 $X$ 是 Noether 概形，每个点都有一个闭点作为其特殊化. 又因正则性在局部化下保持（见定理~\ref{item:ag-3-1-10} v)），$X$ 正则当且仅当它在所有闭点处的局部环都正则. 设 $x$ 是闭点，由命题~\ref{item:ag-3-1-8}和~\ref{item:ag-3-1-9}，闭浸入 $\operatorname{Spec}(\kappa(x))\hookrightarrow X$ 正则当且仅当 $x$ 是正则点.
\textup{ii)} 由定理~\ref{item:ag-3-1-10} iii) 得到.
\end{proof}
注意，切勿将上述陈述中关于 $X$ 和 $Y$ 的假设与结论互换，否则结论就完全不成立了.
\begin{example}\label{item:ag-3-1-13}
设 $X$ 是代数闭域上的簇. 回顾，若概形 $X$ 正则，则称这样的簇 $X$ 为\newterm{ag-smooth-variety}{光滑簇}{smooth variety}. 在此情形下，前一命题~\ref{item:ag-3-1-12}的 iii) 说明，$X$ 的每个光滑子簇 $Y$ 都正则浸入 $X$. 这样的对 $(Y,X)$ 有时称为\newterm{ag-smooth-pair}{光滑对}{smooth pair}. 此外，在这种情形下，纤维 $\mathcal N_{Y/X}\otimes\kappa(x)$ 是 $Y$ 的切空间在 $X$ 于 $x$ 处的余切空间中的正交补.
然而，$Y$ 正则浸入 $X$ 并不意味着 $Y$ 光滑. 例如，$\mathbb P^n$ 中由单个齐次多项式定义的每个超曲面都正则浸入其中. 尽管如此，正则浸入光滑簇的子簇仍具有足够有趣的性质，以致人们给它们起了名字，称之为\newterm{ag-complete-intersection}{完全交}{complete intersection}.
在 $X=\mathbb A_k^n$ 是 $n$ 维仿射空间的情形下，我们来具体说明前一定理的 ii). 设 $x$ 是 $Y$ 的一点，$m:=\dim_x(Y)$. 则 $Y$ 在 $x$ 处光滑当且仅当存在多项式 $g_1,\ldots,g_{n-m}$，它们生成 $Y$ 在 $\mathcal O_{X,x}$ 中的理想，且它们在 $x$ 处的微分（即在 $\mathfrak m_x/\mathfrak m_x^2$ 中的像）线性无关. 换一种观点，假设 $Y$ 由不定元 $T_1,\ldots,T_n$ 的多项式 $f_1,\ldots,f_r$ 定义. 则点 $x\in Y$ 正则当且仅当 Jacobi 矩阵 $\bigl(\frac{\partial f_i}{\partial T_j}(x)\bigr)_{i,j}$ 的秩等于 $n-m$. 事实上，这由如下映射为同构而得：
\begin{align}
\begin{aligned}
\mathfrak m_x/\mathfrak m_x^2&\longrightarrow k^n,\\
f&\longmapsto\left(\frac{\partial f}{\partial T_j}(x)\right)_{j=1,\ldots,n}.
\end{aligned}\tag{3.1.13.1}\label{eq:ag-3-1-13-1}
\end{align}
在 $x$ 处作 Taylor 展开便可看出它确为同构.\unskip\nobreak\hfill\ensuremath{\diamond}
\renewcommand{\envenddiamond}{}
\end{example}
现在给出正则浸入的一些性质.
\begin{proposition}\label{item:ag-3-1-14}
在图~\eqref{eq:ag-3-1-4-1}中，假设浸入 $Y\hookrightarrow X$ 正则. 则在命题~\ref{item:ag-3-1-5}的两个条件之一成立时，浸入 $Y'\hookrightarrow X'$ 也正则.
\end{proposition}
\begin{proof}
归结到局部情形，其中 $X=\operatorname{Spec}(A)$、$Y=\operatorname{Spec}(A/I)$、$X'=\operatorname{Spec}(A')$、$Y'=\operatorname{Spec}(A'/I')$，且 $I'=IA'$. 在命题~\ref{item:ag-3-1-5}的情形 i) 中，同构~\eqref{eq:ag-3-1-1-3}以及对称代数的基变换同构诱导下图中的竖直同构：
\begin{align*}
\begin{tikzcd}[ampersand replacement=\&,column sep=large]
A'/I'\otimes_A\operatorname{Sym}_{A/I}(I/I^2)\arrow[r]\arrow[d,"\sim"'] \& A'/I'\otimes\operatorname{gr}^{I}(A)\arrow[d,"\sim"] \\
\operatorname{Sym}_{A'/I'}(I'/I'^2)\arrow[r] \& \operatorname{gr}^{I'}(A')
\end{tikzcd}
\end{align*}
由此立即得到结论.
现在考虑命题~\ref{item:ag-3-1-5}的情形 ii)，并沿用该命题证明中的记号. 令 $B=A/I$，特别地，有投影 $A\longrightarrow A/I$ 的截面 $B\longrightarrow A$，从而有分解 $A=B\oplus I$. 只须证明，对每个 $n$，态射 $A'\otimes_A I^n=B'\otimes_B I^n\longrightarrow I'^n$ 是同构. 已知它对所有 $n$ 都满射，且在 $n=1$ 时双射. 假设已证明次数 $n$ 时的单射性. 在 $B$-模正合列 $I^{n+1}\hookrightarrow I^n\twoheadrightarrow I^n/I^{n+1}$ 中，右端项在 $B$ 上平坦，因为由假设，它是局部自由 $B$-模 $I/I^2$ 的一个对称幂. 因此，将标量从 $B$ 扩张到 $B'$，或等价地，从 $A$ 扩张到 $A'$ 后，此列仍正合. 应用次数 $n$ 时的归纳假设，得到正合列
\begin{align*}
0\longrightarrow A'\otimes_A I^{n+1}\longrightarrow I'^n\longrightarrow I'^n/I'^{n+1}\longrightarrow0,
\end{align*}
这说明 $A'\otimes_A I^{n+1}=I'^{n+1}$.
\end{proof}
\begin{proposition}\label{item:ag-3-1-15}
设 $Z\xrightarrow jY\xrightarrow iX$ 是两个闭浸入，则 $\mathcal O_Z$-模列
\begin{align*}
0\longrightarrow j^*\mathcal N_{Y/X}\longrightarrow\mathcal N_{Z/X}\longrightarrow\mathcal N_{Z/Y}\longrightarrow0
\end{align*}
在右端及中间正合（这里的态射均为~\eqref{eq:ag-3-1-4-2}中的类型）. 若 $j$ 正则，则此列在左端也正合. 若 $i$ 与 $j$ 都正则，则它们的复合也正则.
\end{proposition}
\begin{proof}
归结到 $X$ 为局部环的谱的情形，设 $Y$ 由理想 $I$ 给出，$Z$ 由包含 $I$ 的理想 $J$ 给出. 陈述中的列于是写成
\begin{align*}
0\longrightarrow A/J\otimes(I/I^2)=I/IJ\longrightarrow J/J^2\longrightarrow(J/I)/(J/I)^2=J/(I+J^2)\longrightarrow0,
\end{align*}
显然在右端及中间正合. 进一步假设 $J/I$ 是 $A/I$ 中的正则理想. 左端态射经由 $I/IJ\longrightarrow I/(I\cap J^2)$ 分解，而 $I/(I\cap J^2)\longrightarrow J/J^2$ 显然是单射，因此只须证明 $I\cap J^2=IJ$. 为此，取 $A$ 中一族元素 $(f_1,\ldots,f_k)$，使其在 $A/I$ 中的像构成生成 $J/I$ 的正则序列. 设 $K$ 为这族元素生成的理想. 考虑态射列
\begin{align*}
A/J[T_1,\ldots,T_k]\longrightarrow A/I\otimes_A\operatorname{gr}^{K}(A)\longrightarrow\operatorname{gr}^{J/I}(A/I).
\end{align*}
第一个箭头是通常的态射 $A/K[T_1,\ldots,T_k]\longrightarrow\operatorname{gr}^{K}(A)$ 与 $A/I$ 的张量积，第二个同构于各箭头 $K^n\longrightarrow(J/I)^n$ 与 $A/J$ 的张量积的直和. 由构造，这两个箭头都是满射，而由假设，它们的复合是双射. 因此二者均为双射. 对每个 $n\in\mathbb N$，由此得到，投影 $A\longrightarrow A/I$ 诱导的映射 $\pi:K^n\longrightarrow(J/I)^n$ 诱导同构
\begin{align*}
K^n/JK^n\xrightarrow\sim(J/I)^n/(J/I)^{n+1}.
\end{align*}
特别地，$\pi$ 的核 $I\cap K^n$ 包含在 $JK^n$ 中. 于是得到 $I\cap K^n=IK^n+(I\cap K^{n+1})$. 由归纳法，对每个 $N>0$ 都有 $I\cap K^n=IK^n+(I\cap K^{n+N})$；再由 $K$-进滤过的分离性，得到 $I\cap K^n=IK^n$. 特别地，取 $n=2$，得到
\begin{align*}
I\cap J^2=I\cap(IJ+K^2)=IJ+(I\cap K^2)=IJ+IK^2=IJ,
\end{align*}
如所欲证.
设 $(f_1,\ldots,f_k)$ 是生成 $I$ 的正则序列，$(f_{k+1},\ldots,f_\ell)$ 是 $J$ 中一族元素，使其在 $A/I$ 中的像构成生成 $J/I$ 的正则序列. 容易验证，序列 $(f_1,\ldots,f_\ell)$ 正则且生成 $J$. 因而，由命题~\ref{item:ag-3-1-8}，若 $i$ 与 $j$ 正则，则复合 $i\circ j$ 也正则. 在此情形下，上述正合列中的所有 $A/J$-模都是自由模.
\end{proof}